\uuid{zsEC}
\titre{Noyau/Image/Théorème du rang/Polynôme}
\auteur{hanani ; mzari}
\organisation{exo7}
\licence{CC BY-NC-SA 4.0}
\datecreate{2020-01-20}
\niveau{L1}
\module{Algèbre}
\chapitre{Application linéaire}
\sousChapitre{Définition}
\difficulte{3}
\theme{Noyau et image, Noyau, Image, Théorème du rang, Polynôme}
\exo7qcmid{lille-2.applications-lineaires.020}

\contenu{
\begin{question}
\qclassification{107.01}
\qtags{motcle=Noyau/Image/Théorème du rang/Polynôme}
On considère l'application linéaire :
$$\begin{array}{rccc}f:&\Rr_2[X]&\to&\Rr_2[X]\\
& P&\to &XP'-X^2P'',  \end{array}$$
où $\Rr_2[X]$ est l'ensemble des polynômes à coefficients réels de degré $\le 2$ et $P'$ (resp. $P''$) est la dérivée première (resp. seconde) de $P$. Quelles sont les assertions vraies ?
\begin{answers}
\bad{$\{1+X^2\}$ est une base de $\ker f $}.
\good{$\{1, X^2\}$ est une base de $\ker f $}.
\bad{$\{1+X\}$ est une base de $\Im f $}.
\good{ $\mbox{rg} (f) =1$}.
\end{answers}
\begin{explanations}  $\ker f=\{P \in \Rr_2[X] \; ; \; XP'-X^2P''=0\}=\{aX^2+b \, ; \; a,b\in \Rr\}$. Donc $\{1,X^2\}$ est une base de $\ker f$. Du théorème du rang, on déduit que $\mbox{rg}(f)=\dim \Im f=1 $ et comme $f(X)=X\neq 0$, $\{X\}$ est une base de $\Im f$.
\end{explanations}
\end{question}
}
